% META: {"id": "TNSI-LANG-CR-013", "chapitre": "TNSI-LANGAGES-ET-PROGRAMMATION", "type_objet": "cours", "capacites_codes": ["C9", "C10"], "sources_inspiration": [], "mode_creation": "ex_nihilo", "status": "generated"}

\section{Paradigmes de programmation}

% ============================================================
% STRATE 1 — Les trois paradigmes sur un meme exemple
% ============================================================

\definition[1]{
Un \textbf{paradigme de programmation} est une façon générale de structurer un programme et de raisonner sur son fonctionnement. Un même langage, et même un même programme, peut mobiliser \textbf{plusieurs paradigmes}.
}

\textbf{Exemple.} Calculer la somme des carrés des nombres pairs d'une liste, dans trois styles différents.

\textbf{Style impératif} (une suite d'instructions qui modifient un état) :
\begin{python}
def somme_carres_pairs_imperatif(liste):
    total = 0
    for x in liste:
        if x % 2 == 0:
            total += x ** 2
    return total
\end{python}

\textbf{Style fonctionnel} (composition de fonctions, sans modifier d'état) :
\begin{python}
def somme_carres_pairs_fonctionnel(liste):
    return sum(x ** 2 for x in liste if x % 2 == 0)
\end{python}

\textbf{Style objet} (les données et les traitements sont regroupés dans une classe) :
\begin{python}
class ListeNombres:
    def __init__(self, valeurs):
        self.valeurs = valeurs

    def somme_carres_pairs(self):
        return sum(x ** 2 for x in self.valeurs if x % 2 == 0)
\end{python}

% BEGIN-VERIFY
% def somme_carres_pairs_imperatif(liste):
%     total = 0
%     for x in liste:
%         if x % 2 == 0:
%             total += x ** 2
%     return total
%
% def somme_carres_pairs_fonctionnel(liste):
%     return sum(x ** 2 for x in liste if x % 2 == 0)
%
% class ListeNombres:
%     def __init__(self, valeurs):
%         self.valeurs = valeurs
%     def somme_carres_pairs(self):
%         return sum(x ** 2 for x in self.valeurs if x % 2 == 0)
%
% donnees = [1, 2, 3, 4, 5, 6]
% resultat_attendu = 2**2 + 4**2 + 6**2
% assert somme_carres_pairs_imperatif(donnees) == resultat_attendu
% assert somme_carres_pairs_fonctionnel(donnees) == resultat_attendu
% assert ListeNombres(donnees).somme_carres_pairs() == resultat_attendu
% END-VERIFY

\margeAppui{
Les trois versions calculent exactement le même résultat : le paradigme change la \textbf{façon d'écrire} le programme, pas ce qu'il calcule (cohérent avec l'idée que la calculabilité ne dépend pas d'un style d'écriture particulier).
}

% ============================================================
% STRATE 2 — Choisir un paradigme
% ============================================================

\propriete[Choisir un paradigme selon le contexte]{
\begin{itemize}
  \item Le style \textbf{impératif} est naturel pour décrire une suite d'étapes avec un état qui évolue (simulation, algorithme pas à pas).
  \item Le style \textbf{fonctionnel} favorise la lisibilité et limite les effets de bord pour des transformations de données (filtrer, transformer, agréger).
  \item Le style \textbf{objet} est adapté quand on manipule des \textbf{entités} avec un état propre et des comportements associés (une classe \lstinline{Compte}, \lstinline{Pile}, \lstinline{Noeud}...).
\end{itemize}
Ce choix n'est presque jamais exclusif : un même projet combine typiquement les trois styles selon les besoins de chaque partie.
}

\erreurFrequente{
\textbf{Croire qu'un langage \og impose \fg{} un seul paradigme.} Python permet d'écrire du code impératif, fonctionnel (fonctions \lstinline{map}, \lstinline{filter}, expressions génératrices) et objet (classes) au sein d'un \textbf{même} programme, voire d'une même fonction. Le paradigme est un choix de conception, pas une contrainte du langage.
}
